\subsection{Free groupoids}

DEFINITION 9. --- Let C be a quiver. The free groupoid constructed on C, denoted \(\mathrm{Grp}(\mathrm{C})\), is the groupoid \(\varpi_{\widetilde{C}}\) of path classes of the graph \(\widetilde{C}\) associated with C.

The set of vertices of Grp(C) is equal to that of C; the orbits of Grp(C) are the connected components of C.

Let C be a nonempty quiver. By Corollary 4, p. 174, the graph \(\widetilde{\mathrm{C}}\) associated with C is a tree if and only if the free groupoid \(\mathrm{Grp}(\mathrm{C})\) constructed on C is simply transitive.

The composite of the canonical quiver morphisms \(i\colon \mathrm{C}\to \widetilde{\mathrm{C}}\) and \(j\colon \widetilde{\mathrm{C}}\to \varpi_{\widetilde{\mathrm{C}}}\) is a quiver morphism from C to \(\mathrm{Grp}(\mathrm{C})\) which we call the canonical morphism from C to \(\mathrm{Grp}(\mathrm{C})\) . Let us denote it \(\theta\) . For every vertex \(a\) of C, we have \(\theta (a) = a\) . If \(f\) is an arrow of C, \(\theta (f)\) is the class of the path \((o(f),(f,1),t(f))\) of \(\widetilde{\mathrm{C}}\) . The morphism \(\theta\) is injective (II, p. 174, cor. 2) and its image generates the groupoid \(\mathrm{Grp}(\mathrm{C})\) .

Let \(\varphi\colon\mathrm{C}\to\mathrm{C}^{\prime}\) be a quiver morphism. Denote by \(\theta_{\mathrm{C}}\) (resp. \(\theta_{\mathrm{C}^{\prime}}\) ) the canonical morphism from C to \(\mathrm{Grp}(\mathrm{C})\) (resp. from \(\mathrm{C}^{\prime}\) to \(\mathrm{Grp}(\mathrm{C}^{\prime})\) ). The morphism \(\varpi(\widetilde{\varphi})\) is the unique groupoid morphism \(\psi\) from \(\mathrm{Grp}(\mathrm{C})\) to \(\mathrm{Grp}(\mathrm{C}^{\prime})\) such that \(\psi\circ\theta_{\mathrm{C}}=\theta_{\mathrm{C}^{\prime}}\circ\varphi\) .

PROPOSITION 6. --- Let C be a quiver, let G be a groupoid, and let \(\varphi\) be a quiver morphism from C to G. There exists a unique groupoid morphism \(\varphi'\) from Grp(C) to G such that \(\varphi' \circ \theta_{\mathrm{C}} = \varphi\) .

Let \(\psi\) be the quiver morphism from \(\widetilde{\mathbf{C}}\) to G such that \(\psi(a) = \varphi(a)\) for every vertex \(a\) of C and \(\psi(f, \varepsilon) = \varphi(f)^{\varepsilon}\) for every arrow \(f\) of C and every \(\varepsilon \in \{-1, 1\}\) . We have \(\psi \circ i = \varphi\) . By Proposition 5, p. 173, there exists a morphism \(\varphi'\) from \(\mathrm{Grp}(\mathrm{C}) = \varpi_{\widetilde{\mathrm{C}}}\) to G such that \(\varphi' \circ j = \psi\) . One then has \(\varphi' \circ \theta_{\mathrm{C}} = \varphi' \circ j \circ i = \psi \circ i = \varphi\) . As in the proof of Prop. 5, the uniqueness of \(\varphi'\) results from the fact that \(\theta_{\mathrm{C}}(\mathrm{C})\) generates the groupoid \(\mathrm{Grp}(\mathrm{C})\) .

Under the hypotheses of Prop. 6, one sometimes says that \(\varphi'\) is the groupoid morphism from Grp(C) to G extending the quiver morphism \(\varphi\) .

Example. --- Let C be a quiver with a single vertex s and let A be the set of its arrows. The groupoid \(\mathrm{Grp}(\mathrm{C})\) is then the groupoid associated with the free group \(\mathrm{F}(\mathrm{A})\) built on A. Indeed, it has a single vertex s, and it follows immediately from Prop. 6 that the canonical map from A to \(\mathrm{Grp}(\mathrm{C})_{s}\) deduced from the canonical morphism from C to \(\mathrm{Grp}(\mathrm{C})\) satisfies the universal property of free groups (A, I, p. 85, prop. 8).

